\documentclass{article}
\usepackage[T1]{fontenc}
\usepackage{lmodern}
\usepackage{graphicx}
\usepackage{amsmath,amssymb}
\usepackage[a4paper,margin=2cm]{geometry}
\usepackage[absolute,overlay]{textpos}
\setlength{\TPHorizModule}{1mm}
\setlength{\TPVertModule}{1mm}
\usepackage[indonesian]{babel}
\usepackage{hyperref}
\setlength{\emergencystretch}{2em}
% unit: Halaman 40

\begin{document}

% === Halaman 40 (python/native) ===

40

• Kita pun dapat membuat variasi cipher transposisi dengan atutan yang kita 

definisikan sendiri.

• Contonya seperti berikut. Misalkan plainteks: ITB GANESHA SEPULUH 

• Bagi menjadi blok-blok 8-huruf. Jika < 8, tambahkan huruf dummy.

• Cipherteks: STBAGNEIUASPEULHGABDCEFH 

1 2 3 4 5 6 7 8 1 2 3 4 5 6 7 8 1 2 3 4 5 6 7 8 

I T B G A N E S H A S E P U L U H A B C D E F G 





S T B A G N E I U A S P E U L H G A B D C E F H 

1 2 3 4 5 6 7 8 1 2 3 4 5 6 7 8 1 2 3 4 5 6 7 8

\end{document}
